\documentclass[11pt]{article}
\usepackage{amsmath,amssymb}
\usepackage[hidelinks]{hyperref}
\title{Exact corridor recovery deadlines after a lost formation observation}
\date{27 November 2026}
\begin{document}
\maketitle

\paragraph{Audit.}
An ordered line of robots starts and finishes at rest. Surviving
robots must preserve declared hard clearances and any required old
geometric links continuously. Additional rank-span density and radio
requirements must hold at the deadline. Robot order, effective
actuator scales, admissible post-fault contexts and the information
available before the maneuver are pinned. There are no subsequent
observations that distinguish those contexts. This is a model of
collision clearance and distance-based radio geometry, not packet
delivery, transient impacts, movable obstacles or a general
controller for robots with uncoordinated clocks. The result is an
R6 synthesis in this family, not a native experiment or a novelty
claim about shortest paths.

\paragraph{What the graph means.}
Give each robot and each of its available local observation cells one
terminal position variable. A shared initial position and actuator
scale are part of that local cell. Keep distinct variables for distinct
cells, and keep fixed landmarks at scale zero. An admitted world
supplies ordinary bounds \(y_j-y_i\le c\) on the variables of its
locally observed cells. Take the union of these directed edges, with
the same variable used wherever the same observation recurs. A lower
span, an upper span, a rank-based density bound, or a no-pass interval
against a fixed obstruction fits this interface. Old hard bounds also
appear as terminal edges and must hold at the initial cut in every
admitted world. The solver works on the whole declared support to
construct a table before motion; no agent is told the actual hidden
world. Each agent reads its own cell in the preinstalled table.

Let \(x_i\) be the cell's initial position and \(k_i\ge0\) its speed
and acceleration scale. Write \(d(i,j)\) for shortest directed
distances in the joined difference graph. A negative cycle rules out
all terminal configurations. Otherwise for a base displacement
budget \(B\) the exact condition is
\[
 x_j-k_jB\le x_i+k_iB+d(i,j)
 \quad\text{for every reachable pair }i,j .
\]
Necessity follows by summing every path and applying the local
travel bounds. For sufficiency choose
\[
 y_j=\min_{i\leadsto j}\{x_i+k_iB+d(i,j)\}.
\]
The self term gives the upper reach bound and the test gives the
lower bound. Appending an edge shows that the vector meets every
constraint, including the bounds shared across worlds. Thus the
minimum \(B\) is the maximum of zero and all positive violations
\((x_j-x_i-d(i,j))/(k_i+k_j)\). A positive numerator with zero
denominator is an impossibility at every time; an unreachable
pair is omitted. These are different obstructions from a finite
budget that simply misses the deadline.

\paragraph{Continuous actuation and hard invariance.}
For single integrators the base travel budget can be \(vT\).
For double integrators fix \(a,v>0\), with each robot satisfying
\(|\ddot x_i|\le k_i a\), \(|\dot x_i|\le k_i v\).
Rest at both ends bounds its velocity magnitude at time \(s\)
by \(k_i\min(as,a(T-s),v)\). Integrating gives \(k_iB(T)\), where
\[
 B(T)=
 \begin{cases}
 aT^2/4,&T\le 2v/a,\\
 v(T-v/a),&T\ge 2v/a .
 \end{cases}
\]
This bound applies to any admissible trajectory, including those
that change direction. Conversely, take the integral of the
nonnegative base envelope and divide it by \(B(T)\) to get
a common progress function \(\lambda(s)\). Use
\(x_i(s)=x_i+\lambda(s)(y_i-x_i)\). These are genuine continuous
rest-to-rest motions (piecewise constant acceleration when
appropriate). Their speeds and accelerations are within every
scaled bound. At every instant this is a convex interpolation
of the initial and terminal positions, so every old hard
difference bound holds throughout. At zero budget the construction
uses the constant path. The deadline is therefore obtained by
inverting the displayed base budget at the graph's exact minimum.
The common acceleration-to-speed ratio and rest cuts matter;
no such construction is asserted for arbitrary different profiles,
moving starts, or nonconvex geometry.

\paragraph{A sharp information cost from hidden density demand.}
Take \(n\) blocks of \(m\ge2\) consecutive gaps, all initially of
size \(g>0\), with unit actuator scales and enough free extent
at the ends. Throughout, require every adjacent gap in \([g,R]\).
Exactly one block has a new maximum-density requirement at the
deadline: its rank-\(m\) span must be at least \(mg+\delta\), with
\(\delta>0\). This is a specified rank-span floor, not a surrogate
claim about arbitrary area coverage. A loss of a separate
diagnostic prevents the agents from learning which block requires
the span; the adjacency requirement remains a geometric constraint.
Availability of that separate message does not follow from the
geometric links.

If \(R<g+\delta/m\), even the identified demand is physically
impossible. Otherwise an observation cell containing \(q\)
possible demanded blocks must widen each of those blocks. Their
disjoint gaps imply a net endpoint span gain of \(q\delta\).
The endpoints can jointly contribute only \(2B(T)\).
Conversely, widen every gap in the identified candidate blocks by
\(\delta/m\), keep the others unchanged, and center the resulting
cumulative displacement to fit in \([-q\delta/2,q\delta/2]\).
This construction preserves every hard gap bound and attains
\(B=q\delta/2\). Knowing the actual block costs only
\(\delta/2\); with no notice the budget needed is \(n\delta/2\).
If a sender knows the demanded block and an identical error-free
message reaches every agent before motion, a message cell can hold
at most \(q_*=\lfloor 2B(T)/\delta\rfloor\) candidates. For
\(q_*\ge1\), exactly \(\lceil n/q_*\rceil\) distinguishable messages
are necessary and sufficient by partitioning the blocks. When
\(q_*=0\), information cannot compensate for inadequate motion.
For example, with \(n=5\), \(\delta=1\), \(a=2\) and a nonbinding
speed limit, the per-world time is 1, the no-message time is
\(\sqrt5\), and at \(T=\sqrt2\) precisely three messages are needed.
This lower bound does not assume a latency-free delivery device
exists. It characterizes what such a delivered pre-motion notice
would have to distinguish.

\paragraph{Proof reuse and comparison.}
A single added edge \(u\to v\) of weight \(c\) creates a negative
cycle exactly when \(d(v,u)+c<0\). Otherwise the revised closure
is the minimum of the old distance and the route
\(d(i,u)+c+d(v,j)\): a best route need not traverse the new
edge twice unless it uses a negative cycle. Joining two previously
distinct observations of the same robot means adding both
directions of a zero-weight equality, checking each update.
The old deadline alone cannot answer this new query. Removing an
edge instead requires recovering any shortest-path alternatives
that used it, or recomputing closure. A tighter actuator scale
without a graph change changes only the corresponding travel
denominators. These computations do not infer a reliable sensor,
scope amendment, or permission to relax hard clearance.

The strongest matched fixed-deadline comparator is the same
deadline's ordinary difference-feasibility problem with all
local indistinguishability equalities and travel bounds; it ties.
The shortest-path reduction uses the classical simple temporal
problem algebra of \href{https://www.sciencedirect.com/science/article/pii/0004370291900066}{Dechter, Meiri and Pearl},
though these variables denote spatial endpoints. It extends the
known-order corridor recovery test by folding across local cells,
interpreting higher-rank spans, and proving the motion and message
cutoffs in the stated subclass.
The \href{https://en.endo-lab.jp/research/detail/499}{Maeda, Endo
and Matsuno experiment} recovers connectivity while avoiding a
failed body; no claim is made to reproduce its controller here.
\href{https://arxiv.org/html/2609.29939v1\#S3.SS2.SSS3}{Managed
make-before-break} instead keeps a departing robot active until
prospective links qualify. It is not a fallback for a missing
pre-motion demand observation. \href{https://arxiv.org/abs/2210.01354}{Breeden and Panagou}
construct joint viability domains under input limits; the
restricted continuous interpolation here supplies an exact
endpoint-to-motion witness instead of just checking that a
single instant's inputs are feasible.

\paragraph{Checks and residual.}
The finite replay checks 22,400 scaled rational feasibility
queries against endpoint enumeration, with 580 reconstructed
attainment witnesses, 645 incremental closures, 1,506 demand
cells and 24 individually physically impossible span controls.
Lean checks eight lemmas, including the real finite-distance
closed-box equivalence; the analytic dynamics and shortest-path
calculation are given above, not represented as a fully formalized
hybrid controller. Actual demand acquisition and delayed or
partially delivered notices are separate obligations. This does
not close R6 or imply an independent mission pass.
\end{document}
